\documentclass[border=4pt]{standalone}
\usepackage{amsmath,amssymb,mathtools,xcolor,varwidth}
\definecolor{radigreen}{HTML}{00A35A}
\begin{document}
\begin{varwidth}{28em}
\large\bfseries
Draw the \textcolor{radigreen}{radii $SA,\,SB,\,SC,\,SD,\,SE,\,SF$} from the centre to every vertex of the path.      As the body advances one step, its radius sweeps across a thin triangle: the area cut out between      successive radii is precisely the area described "by a radius drawn to the centre." Kepler saw this pattern      in the wanderings of Mars — the planet moves fast near the Sun and slow when far, yet always paints equal      areas in equal times. Newton's aim is to \emph{explain} it: he will prove that triangle $SAB$, then $SBC$,      then $SCD$, and so on are all exactly equal, so that the running total of \textcolor{radigreen}{swept area} climbs in      perfect lockstep with the clock. In modern language this equal-area law is nothing but the conservation of      angular momentum, hiding in plain geometric sight.
\end{varwidth}
\end{document}
